% -----------------------------------------------------------------------------
% Differential Geometry lecture notes preamble — ocean atlas style
% Structure: thmtools + mdframed
% Compile with a class that has chapters, e.g. \documentclass{report} or book.
% -----------------------------------------------------------------------------

% Basic commands
\usepackage[margin=1.45in]{geometry}
\usepackage[spanish, es-noquoting]{babel}
\usepackage[T1]{fontenc}
\usepackage[utf8]{inputenc}
\usepackage{lmodern}
\usepackage{microtype}
\usepackage{amsmath, amsfonts, mathtools, amsthm, amssymb}
\usepackage[usenames,dvipsnames,svgnames]{xcolor}
\usepackage{float}
\usepackage{enumitem}
\setlist{itemsep=0.15em, topsep=0.35em}

% TikZ and diagrams
\usepackage{tikz}
\usepackage{tikz-cd}
\usetikzlibrary{intersections, angles, quotes, calc, positioning}
\usetikzlibrary{arrows.meta, decorations.pathmorphing, shadows.blur}
\usepackage{pgfplots}
\pgfplotsset{compat=1.18}

% Hyperlinks
\usepackage[colorlinks=true, linkcolor=black, citecolor=black, urlcolor=black]{hyperref}

% Theorem machinery
\usepackage{thmtools}
\usepackage[framemethod=TikZ]{mdframed}
\mdfsetup{skipabove=1em, skipbelow=0.2em}

% Useful math commands
\newcommand{\R}{\mathbb{R}}
\newcommand{\C}{\mathbb{C}}
\newcommand{\F}{\mathbb{F}}
\newcommand{\N}{\mathbb{N}}
\newcommand{\Q}{\mathbb{Q}}
\newcommand{\Z}{\mathbb{Z}}
\newcommand{\K}{\mathbb{K}}
\newcommand{\eps}{\varepsilon}
\let\epsilon\varepsilon
\newcommand{\ds}{\displaystyle}
\newcommand{\ol}{\overline}
\newcommand{\ul}{\underline}
\newcommand{\xto}[1]{\xrightarrow{#1}}
\newcommand{\xfrom}[1]{\xleftarrow{#1}}
\DeclareMathOperator{\Ker}{Ker}
\DeclareMathOperator{\Imagen}{Im}
\DeclareMathOperator{\Hom}{Hom}
\DeclareMathOperator{\Aut}{Aut}
\DeclareMathOperator{\End}{End}
\DeclareMathOperator{\GL}{GL}
\DeclareMathOperator{\SL}{SL}
\DeclareMathOperator{\SO}{SO}
\DeclareMathOperator{\grad}{grad}
\DeclareMathOperator{\traz}{traz}
\DeclareMathOperator{\Int}{Int}
\DeclareMathOperator{\cl}{cl}
\DeclareMathOperator{\id}{id}
\DeclareMathOperator{\dist}{dist}
\DeclareMathOperator{\rank}{rank}
\DeclareMathOperator{\Span}{span}
\DeclareMathOperator{\diam}{diam}
\DeclareMathOperator{\supp}{supp}
\DeclareMathOperator{\Spec}{Spec}
\DeclareMathOperator{\Var}{Var}
\DeclareMathOperator{\Cov}{Cov}
\DeclareMathOperator{\ord}{ord}
\DeclareMathOperator{\dom}{dom}
\DeclareMathOperator{\ran}{ran}
\DeclareMathOperator{\mcd}{mcd}
\DeclareMathOperator{\mcm}{mcm}

% Header/footer
\usepackage{fancyhdr}
\pagestyle{fancy}
\fancyhf{}
\fancyhead[L]{\sffamily\nouppercase{Victoria Eugenia Torroja}}
\fancyhead[R]{\sffamily\nouppercase\rightmark}
\fancyfoot[C]{\sffamily\leftmark}
\fancyfoot[R]{\sffamily\thepage}
\renewcommand{\headrulewidth}{0.4pt}
\renewcommand{\footrulewidth}{0pt}

% Generic proof name in Spanish
\renewcommand{\proofname}{Demostración}

% Palette: ocean blue and soft seafoam
\definecolor{GeoOcean}{HTML}{063A5B}
\definecolor{GeoTeal}{HTML}{087E8B}
\definecolor{GeoFoam}{HTML}{EEFDFB}
\definecolor{GeoSky}{HTML}{F0F8FF}
\definecolor{GeoSand}{HTML}{F8F3E8}

\definecolor{GeoCoral}{HTML}{C44536}
\definecolor{GeoRose}{HTML}{D95D73}
\definecolor{GeoBlush}{HTML}{FFF0F2}
\definecolor{GeoWine}{HTML}{6D213C}

\colorlet{ResultTitle}{GeoOcean}
\colorlet{DefinitionTitle}{GeoTeal!80!black}
\colorlet{ExampleTitle}{GeoOcean!70}
\colorlet{NotationTitle}{GeoTeal!70!black}

% mdframed base styles
\mdfdefinestyle{georesultframe}{%
nobreak,
    roundcorner=12pt,
    linewidth=0.75pt,
    linecolor=GeoCoral!80!black,
    backgroundcolor=GeoBlush,
    leftmargin=0pt,
    rightmargin=0pt,
    innerleftmargin=10pt,
    innerrightmargin=10pt,
    innertopmargin=8pt,
    innerbottommargin=8pt,
    skipabove=0.95em,
    skipbelow=0.95em,
    shadow=true,
    shadowsize=4pt,
    shadowcolor=black!18,
    splittopskip=1.2em,
    splitbottomskip=0.8em,
    %firstextra={\draw[GeoOcean, line width=3.5pt] (P) -- (P|-O);\draw[GeoTeal!60, line width=0.7pt] (O) -- (O|-P);},
    %middleextra={\draw[GeoOcean, line width=3.5pt] (P) -- (P|-O);\draw[GeoTeal!60, line width=0.7pt] (O) -- (O|-P);},
    %secondextra={\draw[GeoOcean, line width=3.5pt] (P) -- (P|-O);\draw[GeoTeal!60, line width=0.7pt] (O) -- (O|-P);},
    %singleextra={\draw[GeoOcean, line width=3.5pt] (P) -- (P|-O);\draw[GeoTeal!60, line width=0.7pt] (O) -- (O|-P);}
}

\mdfdefinestyle{geodefinitionframe}{%
nobreak,
    roundcorner=12pt,
    linewidth=0.75pt,
    linecolor=GeoTeal,
    backgroundcolor=GeoFoam,
    leftmargin=0pt,
    rightmargin=0pt,
    innerleftmargin=10pt,
    innerrightmargin=10pt,
    innertopmargin=8pt,
    innerbottommargin=8pt,
    skipabove=0.95em,
    skipbelow=0.95em,
    shadow=true,
    shadowsize=4pt,
    shadowcolor=black!18,
    splittopskip=1.2em,
    splitbottomskip=0.8em,
    %firstextra={\draw[GeoTeal, line width=3.5pt] (P) -- (P|-O);},
    %middleextra={\draw[GeoTeal, line width=3.5pt] (P) -- (P|-O);},
    %secondextra={\draw[GeoTeal, line width=3.5pt] (P) -- (P|-O);},
    %singleextra={\draw[GeoTeal, line width=3.5pt] (P) -- (P|-O);}
}

\mdfdefinestyle{geoexampleframe}{%
    roundcorner=7pt,
    linewidth=0.75pt,
    linecolor=GeoOcean!50,
    backgroundcolor=GeoSand,
    leftmargin=0pt,
    rightmargin=0pt,
    innerleftmargin=10pt,
    innerrightmargin=10pt,
    innertopmargin=8pt,
    innerbottommargin=8pt,
    skipabove=0.95em,
    skipbelow=0.95em,
    shadow=true,
    shadowsize=3pt,
    shadowcolor=black!10,
    splittopskip=1.2em,
    splitbottomskip=0.8em
}

\mdfdefinestyle{geonotationframe}{%
    roundcorner=7pt,
    linewidth=0.75pt,
    linecolor=GeoTeal!60!black,
    backgroundcolor=white,
    leftmargin=0pt,
    rightmargin=0pt,
    innerleftmargin=10pt,
    innerrightmargin=10pt,
    innertopmargin=8pt,
    innerbottommargin=8pt,
    skipabove=0.95em,
    skipbelow=0.95em,
    shadow=true,
    shadowsize=3pt,
    shadowcolor=black!10,
    splittopskip=1.2em,
    splitbottomskip=0.8em
}

% Differential geometry helpers
\DeclareMathOperator{\Ric}{Ric}
\DeclareMathOperator{\Scal}{Scal}
\DeclareMathOperator{\vol}{vol}
\DeclareMathOperator{\Lie}{Lie}
\newcommand{\dd}{\,\mathrm{d}}
\newcommand{\LieD}{\mathcal{L}}
\newcommand{\T}{\mathrm{T}}

% Theorem styles
\declaretheoremstyle[
    headfont=\bfseries\sffamily\color{DefinitionTitle},
    bodyfont=\normalfont,
    spaceabove=0pt,
    spacebelow=0pt,
    mdframed={style=geodefinitionframe}
]{geodefinitionbox}

\declaretheoremstyle[
    headfont=\bfseries\sffamily\color{GeoCoral!80!black},
    bodyfont=\normalfont,
    spaceabove=0pt,
    spacebelow=0pt,
    mdframed={style=georesultframe}
]{georesultbox}

\declaretheoremstyle[
    headfont=\bfseries\sffamily\color{Rhodamine!50!black},
    bodyfont=\normalfont,
    spaceabove=0pt,
    spacebelow=0pt,
    mdframed={nobreak,
    roundcorner=12pt,
    linewidth=0.75pt,
    linecolor=RoyalPurple!20!black,
    backgroundcolor=Rhodamine!10,
    leftmargin=0pt,
    rightmargin=0pt,
    innerleftmargin=10pt,
    innerrightmargin=10pt,
    innertopmargin=8pt,
    innerbottommargin=8pt,
    skipabove=0.95em,
    skipbelow=0.95em,
    shadow=true,
    shadowsize=4pt,
    shadowcolor=black!18,
    splittopskip=1.2em,
    splitbottomskip=0.8em,}
]{georesultbox1}

\declaretheoremstyle[
    headfont=\bfseries\sffamily\color{ExampleTitle},
    bodyfont=\normalfont,
    spaceabove=0pt,
    spacebelow=0pt,
    mdframed={style=geoexampleframe}
]{geoexamplebox}

\declaretheoremstyle[
    headfont=\bfseries\sffamily\color{ExampleTitle},
    bodyfont=\normalfont,
    spaceabove=0pt,
    spacebelow=0pt,
    mdframed={
    linecolor = CadetBlue!80!black,
    linewidth = 2pt,
    topline = false,
    bottomline = false, 
    rightline = false,
    }
]{geoobsbox}

\declaretheoremstyle[
    headfont=\bfseries\sffamily\color{NotationTitle},
    bodyfont=\normalfont,
    spaceabove=0pt,
    spacebelow=0pt,
    mdframed={style=geonotationframe}
]{geonotationbox}

% Environments. Same names as your earlier notes.
\declaretheorem[style=geodefinitionbox, name=Definición, numberwithin=chapter]{definition}
\declaretheorem[style=georesultbox, name=Teorema, numberwithin=chapter]{theorem}
\declaretheorem[style=georesultbox, name=Proposición, numberwithin=chapter]{prop}
\declaretheorem[style=georesultbox, name=Lema, numberwithin=chapter]{lema}
\declaretheorem[style=georesultbox, name=Corolario, numberwithin=chapter]{colorary}
\declaretheorem[style=geoobsbox, numbered=no, name=Ejemplo]{eg}
\declaretheorem[style=geoexamplebox, numbered=no, name=Observación]{observation}
\declaretheorem[style=geonotationbox, numbered=no, name=Notación]{notation}
\declaretheorem[style=geoexamplebox, numbered=no, name=Comentario]{remark}

